\documentclass[10pt,a4paper]{amsart}
\usepackage[margin=19mm]{geometry}
\usepackage{amsmath,amssymb,mathtools}
\usepackage[scr=rsfs,cal=euler]{mathalfa}
\usepackage{xfrac}
\usepackage[unicode,hidelinks,bookmarks=false]{hyperref}
\newcommand{\ie}{i.e.\ }
\newcommand{\cf}{cf.\ }
\newcommand{\resp}{respectively\ }
\newcommand{\Subsection}{\subsection}
\providecommand{\nobreakdash}{\nobreak\hbox{-}}
\setlength{\emergencystretch}{2em}
\multlinegap0pt
\usepackage{bm}
\usepackage{bbm}
\usepackage{dsfont}
\usepackage{xspace}
\usepackage{cancel}
\usepackage{mathtools}
\usepackage{amsthm}
\makeatletter
\@ifundefined{theorem}{\newtheorem{theorem}{Theorem}}{}
\makeatother
\title[Boundary value problems for some quasilinear parabolic equat]{Boundary value problems for some quasilinear parabolic equations under minimal conditions on the coefficients}
\author{S.G. Pyatkov}
\date{Source: arXiv:2606.16446v1 (CC BY 4.0)}
\begin{document}
\maketitle

\noindent\emph{Source and licence.} Boundary value problems for some quasilinear parabolic equations under minimal conditions on the coefficients. Version arXiv:2606.16446v1 (CC BY 4.0), distributed under the Creative Commons Attribution 4.0 International licence, as recorded in the arXiv licence metadata for this exact version and shown on the abstract page https://arxiv.org/abs/2606.16446v1. Full rights notice: licenses/arxiv-2606.16446-CC-BY-4.0.txt.\par\smallskip

\noindent\textit{Editorial scope.} Counted unit: the Theorem whose source label is \texttt{thr3} in the article (source0.tex lines 388--394, source label \texttt{thr3}), delivered together with the article's own complete original proof of it (source0.tex lines 395--461, one \texttt{proof} environment of 67 lines). No mathematics is written, repaired or re-derived: statement and proof are the article's own text line for line. Required context retained uncounted: e1 (lines 74--76); e2 (lines 79--81); e3 (lines 83--85); e6 (lines 229--231); e8 (lines 239--241); thr2 (lines 265--270); e11 (lines 287--291); e15 (lines 312--315); e18 (lines 346--349); e201 (lines 378--380). Every definition, display and label of those lines is retained unchanged. Omitted from this package and not redistributed: 1--73, 77--78, 82--82, 86--228, 232--238, 242--264, 271--286, 292--311, 316--345, 350--377, 381--387, 462--1418. Nothing inside a retained excerpt is omitted or altered: every retained body is delivered exactly as the article's lines are written, so no omission receipt of any kind accompanies this package. Citations: the retained excerpts carry no citation command at all, so no bibliography entry of the article is transcribed and no reference is redistributed. Reference adaptation: none was needed and none was made; every reference inside the delivered unit resolves inside it, so no unresolved reference or citation is produced. Printed slips kept literal: no printed word, symbol, hypothesis or number of the counted statement or of its proof was altered. Layout and class: the article's own class is \texttt{article} with packages including fontenc, inputenc, babel, amsmath, amsfonts, amssymb, graphicx, amsthm, geometry; the wrapper is the lane's pinned \texttt{amsart} editorial wrapper at 10pt on A4, which restates the theorem-like environments the retained text needs and adds the article's own macro definitions that the retained text needs (none), with extra wrapper packages (bm, bbm, dsfont, xspace, cancel, mathtools). The delivered numbering is the unit's own. Assets: no retained span contains a figure environment, a graphics-inclusion command, a PDF-page inclusion, a table or a TikZ picture; no image file is copied into this package, the package declares no asset dependency and no image file is redistributed with it. Licence: version arXiv:2606.16446v1 of this article is distributed under the Creative Commons Attribution 4.0 International licence, as recorded in the arXiv licence metadata for this exact version and shown on the abstract page https://arxiv.org/abs/2606.16446v1. A full rights notice is delivered at licenses/arxiv-2606.16446-CC-BY-4.0.txt. Characters: every retained excerpt is pure ASCII, so the delivered engine, XeLaTeX with its default Latin Modern text font, reports no missing character.
\begin{equation}\label{e1}
	Mu= u_t-\sum_{i,j=1}^n \partial_{x_i}(a_{ij}(t,x,u)u_{x_j})+ b(x,t,u,\nabla u)=0,
\end{equation}

\begin{equation}\label{e2}
	u|_{t=0}=u_0,\ \ u|_{S}= g(t,x), 
\end{equation}

\begin{equation}\label{e3}
	u|_{t=0}=u_0,\ \ \frac{\partial u}{\partial N}+ \psi(t,x,u)|_{S}=0, \ \frac{\partial u}{\partial N}= \sum_{i,j=1}^na_{ij}u_{x_j}\nu_i.
\end{equation}

\begin{equation}\label{e6}
\sum_{i,j=1}^na_{ij}p_ip_j\geq \delta_0|p|^2\  \forall \vec{p}=(p_1,\ldots,p_n) \in {\mathbb R}^n, \ (t,x)\in Q, \ u\in {\mathbb R},
\end{equation}

\begin{equation}\label{e8}
|b(t,x,u, \vec{p})|\leq \beta_1 |\vec{p}|^2 +g_3(t,x)\   \forall  (t,x)\in Q, \  u\in [-R,R], \ \vec{p}\in {\mathbb R}^n,
\end{equation}

\begin{theorem} \label{thr2}
Assume that  $u\in W_p^{1,2}(Q)$ $(p>(n+2)/2)$ is a  solution to the problem  {\rm \eqref{e1}-\eqref{e2}}, the conditions 
   {\rm \eqref{e6}-\eqref{e8}} hold, and   $g(t,x)\in L_\infty(S)$, $u_0\in L_\infty(G)$. 
Then there exists a constant  $M>0$, depending on $\delta_0, \beta_0$,  $\|g_2\|_{L_{p}(Q)}$, $\|u\|_{L_\infty(\Gamma_T)}$
such that  $\|u\|_{L_\infty(Q)}\leq M$. 
    \end{theorem}

\begin{multline}\label{e11}
\frac{1}{m}\frac{\partial}{\partial t} \int_G (w^{(k)})^2\,dx +\int_{A_k(t)} 
\sum_{i,j=1}^n a_{ij}v_{x_j}v_{x_i} [(m-1)|v|^{m-2}w^{(k)}+ m|v|^{2m-2}] \\ + 
e^{-\lambda t}b(t,x,ve^{\lambda t},e^{\lambda t}\nabla v)|v|^{m-2} v w^{(k)}\,dx=0.
\end{multline}

\begin{equation}\label{e15}
\max_{t\in (0,T)} \int_{G} (w^{(k)})^2\,dx +
\int_{Q} |\nabla w^{(k)}|^2+ \lambda  w w^{(k)}\,dxdt \leq  c_0\int_{L_k} f w w^{(k)} \,dxdt, 
\end{equation}

 \begin{equation}\label{e18}
\max_{t\in (0,T)} \int_{G} (w^{(k)})^2\,dx +
\int_{Q} |\nabla w^{(k)}|^2+ \lambda  (w^{(k)})^2\,dx \leq  2c_2k^2 (\mu(L_k))^{1/p'}.
\end{equation}

\begin{equation}\label{e201}
\psi(t,x,u)sgn\,u\geq -g_4-\beta_3 |u|, 
\end{equation}

\begin{theorem} \label{thr3}
 Assume that  $u\in W_p^{1,2}(Q)$ $(p>(n+2)/2)$ is a solution to the problem  {\rm \eqref{e1}, \eqref{e3}}, the conditions 
   {\rm \eqref{e6}-\eqref{e8}, \eqref{e201}} hold, and    $u_0\in L_\infty(G)$. 
Then there exists a constant  $M>0$, depending on  $\beta_0$,  $\|g_2\|_{L_{p}(Q)}$,  $\|u_0\|_{L_\infty(G)}$, and
$\|g_4\|_{W_p^{s_0,2s_0}(S)}$ such that 
 $\|u\|_{L_\infty(Q)}\leq M$. 
    \end{theorem}

\begin{proof}
The proof is in line with that of Theorem \ref{thr2}. Repeating the arguments, we arrive at \eqref{e11} with an additional  summand $J= -\int_{\Gamma}e^{-\lambda t} \psi(t,x,e^{\lambda t}v) |v|^{m-2}v w^{(k)}\,d\Gamma$  on the right-hand side. Estimate it.   
In view of the conditions on  $\psi$, we have 
$$
J\leq \int_\Gamma (g_4e^{-\lambda t}+\beta_3|v|)|v|^{m-1}w^{(k)}\,d\Gamma.
$$
 Next, 
\begin{equation*}
\int_G \Delta\rho \beta_3|v| |v|^{m-1}w^{(k)}\,dx + \int_G \nabla \rho \cdot\nabla \beta_3|v|^{m}w^{(k)}\,dx= \int_{\Gamma}\nabla \rho\cdot \nu 
\beta_3|v|^{m}w^{(k)}\,d\Gamma. 
\end{equation*}
Therefore, we derive the inequality 
\begin{multline*}
 \int_{\Gamma} \beta_3|v|^{m}w^{(k)}\,d\Gamma \leq 
 2\int_{\Gamma}\nabla \rho\cdot \nu \beta_3|v|^{m}w^{(k)}\,d\Gamma\leq  \\
 2\int_G \Delta\rho \beta_3|v|^{m}w^{(k)}\,dx+ 2\int_G \nabla \rho \cdot\nabla (\beta_3|v|^{m}w^{(k)})\,dx\leq \\
c_1 \int_G \beta_3|v|^{m}w^{(k)}\,dx+ 2\int_G |\nabla (\beta_3|v|^{m}w^{(k)})|\,dx.
\end{multline*}
The first summand here is estimated by  the quantity
\begin{equation}\label{w21}
\int_G  \beta_3|v|^{m}w^{(k)}\,dx\leq  \beta_3\int_G w w^{(k)}\,dx\leq c_5\int_G (w^{(k)})^2\,dx+ c_5k^2 \mu(A_k(t)). 
\end{equation}
The estimate for the second summand  is of the form 
\begin{equation}\label{w22}
 \int_G |\nabla \beta_3|v|^{m}w^{(k)}|\,dx\leq  \varepsilon \int_{G} |\nabla w^{(k)}|^2  + c(\varepsilon) ((w^{(k)})^2\,dx +k^2 \mu(A_k(t))).
\end{equation}
Next, we use the equality 
\begin{equation*}
\int_G \Delta\rho_1  |v|^{m-2}v w^{(k)}\,dx+ \int_G \nabla \rho_1 \cdot\nabla (|v|^{m-2}v w^{(k)})\,dx= \int_{\Gamma} g_4 |v|^{m-2} v w^{(k)}\,d\Gamma 
\end{equation*}
which guarantees the relation
\begin{equation*}
 |\int_{\Gamma}e^{-\lambda t} g_4 |v|^{m-2} v w^{(k)}\,d\Gamma|  \leq 
 \int_{\Gamma}|\nabla \rho_1| |\nabla ( |v|^{m-2}v w^{(k)})|\,d\Gamma + \frac{1}{m_0}\int_G |\Delta\rho_1|  |v|^{m} w^{(k)}\,dx.
 \end{equation*}
The former  summand on the right-hand side of the previous inequality is bounded  by 
\begin{equation}\label{w24}
 \int_{G}|\nabla \rho_1| |\nabla |v|^{m-2}v w^{(k)}|\,dx\leq  \varepsilon \int_{G} |\nabla w^{(k)}|^2  + c(\varepsilon) (|\nabla \rho_1|^2 (w^{(k)})^2 + k^2\chi_{A_k(t)})\,dx.
\end{equation}
where $\chi_{A_k(t)}$ is the indicator of the set  $A_k(t)$.
Thus, the estimates   \eqref{w24} validate the relation
$$
\int_0^T \int_Ge^{-\lambda t} g_4 |v|^{m-2} v w^{(k)}\,dxdt \leq \varepsilon \int_{Q}|\nabla w^{(k)}|^2 + c(\varepsilon) \int_Q\tilde{g}_3( (w^{(k)})^2 +k^2\chi_{L_k}(t,x))\,dxdt,
$$
 where  $\tilde{g}_3=|\Delta \rho_1|+|\nabla \rho_1|^2$.
By the embedding theorems,  $\tilde{g}_3\in L_p(Q)$.  
The same arguments as those involved in the estimate of the right-hand side 
in \eqref{e15},  ensure the inequality
 $$
 \int_Q\tilde{g}_3 (w^{(k)})^2dxdt\leq \varepsilon_1 |w^{(k)}|_1^2  + c(\varepsilon_1) \int_Q |w^{(k)}|^2\,dxdt, 
$$
with $\varepsilon_1$ an arbitrary positive constant.
Moreover, 
 $$
 k^2 \int_Q \tilde{g}_3 \chi_{L_k(t,x)} \,dxdt\leq k^2\|\tilde{g}_3\|_{L_p(Q)}(\mu(L_k))^{1/p'},\  p'=p/(p-1). 
$$ 
These inequalities  for an appropriate  $\varepsilon_1$  and \eqref{w21}, \eqref{w22} validate the estimate 
$$
\int_0^T J\leq 3\varepsilon \int_{Q} |\nabla w^{(k)}|^2 + c(\varepsilon)( \int_Q |w^{(k)}|^2\,dxdt + k^2 (\mu(L_k))^{1/p'}).
$$
Thus, the inequality  \eqref{e18} from the previous theorem for an appropriate  $\lambda$ and $\varepsilon$ is replaced with 
 \begin{equation}\label{w25}
\max_{t\in (0,T)} \int_{G} (w^{(k)})^2\,dx +
\int_{Q} |\nabla w^{(k)}|^2+ \lambda  (w^{(k)})^2\,dx \leq  c_4 k^2 (\mu(L_k))^{1/p'}.
\end{equation} 
     Further, we can repeat the arguments of the previous theorem.  
\end{proof} 

\end{document}
